\SourceLeaf{158}Bestimmung der Polhöhe aus $p$ unteren und $q$ oberen Kulminationen das Gewicht $\dfrac{4pq}{p+q}$ beigelegt wird, was auch ganz richtig ist. Die Polhöhe ist aber nicht die Summe, sondern nur die \textls{halbe} Summe der Zenitdistanzen bei oberer und unterer Kulmination. Der Fall ist also ganz verschieden. Allgemein lässt sich \textls{ohne viel Mühe} aus den Grundsätzen meiner Theoria und \textls{äußerst leicht} mit meiner neuen Begründung zeigen, dass, --- wenn eine Größe $w$ aus der Kombination der Größen $x,y,z$ etc. entsteht, sodass $w=\alpha x+\beta y+\gamma z+\text{etc.}$, und wenn ferner die Werte von $x,y,z$ etc. unabhängig voneinander so ermittelt sind, dass diesen Bestimmungen die jeweiligen Gewichte $p,q,r$ etc. beizulegen sind, --- der daraus hervorgehende Wert von $w$ eine Zuverlässigkeit hat, deren Gewicht durch $\frac{1}{\dfrac{\alpha\alpha}{p}+\dfrac{\beta\beta}{q}+\dfrac{\gamma\gamma}{r}+\text{etc.}} $ ausgedrückt wird.

(Das Gewicht ist übrigens immer dem Quadrat der Genauigkeit direkt oder dem Quadrat des sogenannten wahrscheinlichen Fehlers umgekehrt proportional. Das ist kein Lehrsatz, sondern bloß die \textls{Definition} des Wortes Gewicht.)



\Needspace{6\baselineskip}
\begin{center}
\textsc{Gauss} an \textsc{Olbers}. Göttingen, 25. Februar 1825.
\end{center}

\begin{center}
. . . . . . . . . . . . . . . . .
\end{center}

P. S. Die Verteilung der Fehler in Gruppen, um die Summe der Quadrate zu finden, bewirkt \textls{jedes Mal}, dass die Summe zu klein wird und die Genauigkeit deshalb größer erscheint, als sie ist.

Die Fehler seien
\[
\begin{array}{lll}
a,\ a',\ a''\ .\ .\ .\ .\ ., & \text{Anzahl}=\alpha, & \text{Mittel}=A, \\
b,\ b',\ b''\ .\ .\ .\ .\ ., & \text{»}\ \ \beta, & \text{»}\ \ B, \\
c,\ c',\ c''\ .\ .\ .\ .\ ., & \text{»}\ \ \gamma, & \text{»}\ \ C, \\
& \text{etc.,} &
\end{array}
\]
Dann gilt genau
\[
aa+a'a'+a''a''\ldots=\alpha A^2+(a-A)^2+(a'-A)^2+(a''-A)^2+\text{etc.}
\]
\[
bb+b'b'+b''b''\ldots=\beta B^2+(b-B)^2+(b'-B)^2+(b''-B)^2+\text{etc.}
\]
\[
cc+c'c'+c''c''\ldots=\gamma C^2+(c-C)^2+(c'-C)^2+(c''-C)^2+\text{etc.}
\]
\[
\text{etc.}
\]

